\documentclass[11pt,a4paper]{report}
\usepackage[utf8]{inputenc}
\usepackage[T1]{fontenc}
\usepackage[french]{babel}
\usepackage{geometry}
\usepackage{amsmath}
\usepackage{amssymb}
\usepackage{graphicx}
\usepackage{xcolor}
\usepackage{booktabs}
\usepackage{array}
\usepackage{multirow}
\usepackage{enumitem}
\usepackage{hyperref}
\usepackage{float}
\usepackage{longtable}
\usepackage{listings}
\usepackage{tikz}
\usetikzlibrary{arrows.meta,positioning,shapes.geometric,calc}
\geometry{margin=2.15cm}
\hypersetup{colorlinks=true,linkcolor=blue,urlcolor=blue}
\definecolor{hardtecblue}{HTML}{145DA0}
\definecolor{hardtecdark}{HTML}{17324D}
\definecolor{hardteclight}{HTML}{EAF3FB}
\definecolor{softgray}{HTML}{F5F7F9}
\definecolor{greenok}{HTML}{0F766E}
\definecolor{warningorange}{HTML}{D97706}
\definecolor{redalert}{HTML}{B91C1C}
\lstset{
    basicstyle=\ttfamily\small,
    backgroundcolor=\color{softgray},
    frame=single,
    breaklines=true,
    numbers=left,
    numberstyle=\tiny,
    showstringspaces=false,
    columns=fullflexible
}
\newcommand{\screenshot}[3]{%
    \begin{figure}[H]
    \centering
    \IfFileExists{rapport_pfa/screenshots/#1}{%
        \includegraphics[width=0.92\textwidth]{rapport_pfa/screenshots/#1}%
    }{%
        \fbox{\begin{minipage}[c][3.2cm][c]{0.88\textwidth}
        \centering\textbf{Screenshot à ajouter}\\[0.2cm]
        \texttt{rapport_pfa/screenshots/#1}\\[0.2cm]
        #2
        \end{minipage}}%
    }
    \caption{#3}
    \end{figure}%
}
\setlength{\parskip}{2pt}
\setlength{\parindent}{1em}
\setlist{nosep,leftmargin=1.25em}
\setlength{\abovecaptionskip}{3pt}
\setlength{\belowcaptionskip}{3pt}
\setlength{\textfloatsep}{8pt plus 2pt minus 2pt}
\setlength{\intextsep}{8pt plus 2pt minus 2pt}
\tikzset{
    component/.style={draw=hardtecdark,fill=hardteclight,rounded corners,minimum width=3.1cm,minimum height=1cm,align=center,font=\small\bfseries},
    data/.style={draw=hardtecdark,fill=green!10,cylinder,shape border rotate=90,aspect=.25,minimum width=2.7cm,minimum height=1.1cm,align=center,font=\small\bfseries},
    decision/.style={draw=hardtecdark,fill=orange!10,diamond,aspect=2,align=center,font=\small\bfseries},
    arrow/.style={-Latex,thick,draw=hardtecdark}
}
\renewcommand{\contentsname}{Table des matières}
\renewcommand{\listfigurename}{Liste des figures}
\renewcommand{\chaptername}{Chapitre}
\title{\textbf{HARDTEC Intelligent Support \& AIOps}\\[0.5cm]\large Plateforme intelligente de support IT, prédiction des incidents et agent IA}
\author{\textbf{Rapport de stage PFA}\\[0.3cm] ENSA Berrechid\\[0.3cm] au sein de HARDTEC MAROC}
\date{2025-2026}

\begin{document}

\maketitle
\tableofcontents
\listoffigures

\chapter*{Remerciements}
\addcontentsline{toc}{chapter}{Remerciements}
Je tiens à remercier profondément Monsieur \textbf{FADIL BADR}, encadrant professionnel chez HARDTEC MAROC, pour son accompagnement, ses conseils avisés et sa confiance. Je remercie également Monsieur \textbf{Pr. NAFIDI Ahmed}, encadrant académique, pour son suivi scientifique et ses orientations pédagogiques. Je remercie aussi le jury composé de \textbf{Pr. JARRAR Abdssamad}, \textbf{Pr. MOUMOUN}, et \textbf{Pr. AHMED NAFIDI}, pour l’intérêt porté à ce travail et pour le temps qu’ils consacreront à son évaluation. Enfin, je remercie toute l’équipe HARDTEC pour son accueil chaleureux et son environnement stimulant.

\chapter*{Résumé}
\addcontentsline{toc}{chapter}{Résumé}
Ce projet de fin d’année porte sur la conception et le développement d’une plateforme intelligente de support IT et d’AIOps au sein de HARDTEC MAROC. La solution ne se limite pas au ticketing : elle associe la classification des demandes, la prédiction du risque d’incident, l’analyse de métriques, la recherche sémantique et un agent IA d’aide à la décision. L’objectif est de rapprocher le support utilisateur et l’observabilité technique dans une même architecture. Le système exploite les données historiques, les modèles de machine learning et les signaux temporels afin de proposer des décisions contextualisées, tout en conservant une validation humaine. Il est déployable de manière reproductible avec Docker et contrôlé par un pipeline CI/CD.

\chapter{Introduction}
\section{Contexte du stage}
Le stage de projet de fin d’année s’inscrit dans le cadre de la formation d’ingénierie à l’ENSA Berrechid. Il a été réalisé au sein de HARDTEC MAROC, une structure active dans les domaines de l’intelligence artificielle, du support IT, des données et des systèmes décisionnels. L’entreprise est confrontée à la gestion d’un volume important de tickets informatiques, souvent hétérogènes, non standardisés et soumis à des délais d’intervention critiques. Les demandes de support portent sur des incidents réseau, problèmes logiciels, incidents matériels, dysfonctionnements applicatifs, demandes de maintenance et cas d’assistance complexes.

Les méthodes traditionnelles de traitement reposent sur une intervention humaine directe, une classification manuelle et un routage parfois subjectif. Cette approche pose plusieurs problèmes :
\begin{itemize}[leftmargin=1.3em]
    \item lenteur de traitement ;
    \item erreurs de tri et d’affectation ;
    \item manque d’historique exploitable ;
    \item coût opérationnel élevé ;
    \item faible capacité à prévoir des pics d’incidents ;
    \item difficulté à capitaliser les meilleures pratiques de résolution.
\end{itemize}

\section{Problématique}
Comment concevoir une plateforme intelligente capable de classer automatiquement les tickets IT, de prédire leur priorité, de les affecter à la bonne file de traitement, de recommander des actions pertinentes et d’exploiter les cas historiques pour améliorer la décision humaine ? Il s’agit de transformer un support réactif en un environnement proactif et assisté par l’IA.

\section{Objectifs du projet}
Le projet a pour objectif principal de mettre en place une solution de support intelligent pour l’entreprise HARDTEC, basée sur les technologies de l’IA, des données et de l’automatisation. Les objectifs spécifiques sont :
\begin{itemize}[leftmargin=1.3em]
    \item développer un système de classification automatique des tickets ;
    \item prédire la priorité et la file d’affectation des incidents ;
    \item exploiter les tickets historiques à travers un moteur RAG ;
    \item détecter le niveau de risque d’incident avec une logique AIOps ;
    \item mettre en place un agent IA pour assister l’opérateur de support ;
    \item déployer le système via Docker et automatiser sa validation avec CI/CD ;
    \item assurer la traçabilité avec MLflow ;
    \item produire un système exploitable dans un contexte professionnel réel.
\end{itemize}

\section{Méthodologie adoptée}
La démarche a suivi une logique progressive :
\begin{enumerate}[leftmargin=1.8em]
    \item analyse du besoin métier et des flux de support ;
    \item étude de l’architecture logicielle existante ;
    \item identification des composants IA et de données ;
    \item conception d’une architecture modulaire ;
    \item développement des modules de prédiction, RAG et AIOps ;
    \item intégration de l’agent IA ;
    \item validation par tests et démonstration ;
    \item préparation du rapport technique et de la synthèse de stage.
\end{enumerate}

\chapter{Présentation de l’entreprise et du contexte professionnel}
\section{Présentation de HARDTEC MAROC}
HARDTEC MAROC se présente publiquement comme un compagnon IT de confiance et propose des solutions digitales destinées à accompagner les entreprises dans leur transformation. Son site officiel, \url{https://hardtec.net}, met en avant les services web, l’accompagnement informatique et la maintenance. Le stage s’inscrit dans ce contexte professionnel en explorant une plateforme d’intelligence artificielle capable de rapprocher le support IT, l’analyse des incidents et l’automatisation de la décision.

\section{Environnement de travail}
Le projet a été développé dans un environnement de travail basé sur le langage Python, les outils d’analyse de données, les modèles d’apprentissage automatique et les services web. Le développement a nécessité l’utilisation de multiples outils :
\begin{itemize}[leftmargin=1.3em]
    \item Python 3.13 ;
    \item FastAPI pour les API ;
    \item Streamlit pour le tableau de bord ;
    \item scikit-learn pour les modèles supervisés ;
    \item FAISS et SentenceTransformers pour le RAG ;
    \item MLflow pour le suivi des expériences ;
    \item SQLite pour le stockage local ;
    \item Docker pour la conteneurisation ;
    \item GitHub Actions pour l’automatisation CI/CD.
\end{itemize}

\section{Synthèse des technologies utilisées}
\begin{longtable}{p{3.5cm}p{5.3cm}p{5.0cm}}
\toprule
\textbf{Domaine} & \textbf{Technologies} & \textbf{Rôle dans le projet} \\
\midrule
Langage & Python 3.13 & Services, traitements, scripts et modèles \\
API & FastAPI, Uvicorn, Pydantic & Services REST et validation des requêtes \\
Interface & Streamlit, Plotly & Tableau de bord et visualisation \\
Machine learning & scikit-learn, NumPy, pandas, SciPy, joblib & Préparation, entraînement et inférence \\
Recherche sémantique & FAISS, SentenceTransformers & Recherche de tickets historiques similaires \\
AIOps & MicroSS, analyse temporelle, détection d’anomalies & Évaluation du risque d’incident \\
Données & SQLite, DuckDB, Snowflake, CSV & Stockage, analytique et échanges de données \\
Transformation & dbt Core, dbt Snowflake & Modèles Bronze, Silver et Gold \\
Orchestration & Apache Airflow & Planification des pipelines batch \\
MLOps & MLflow & Runs, métriques et artefacts de modèles \\
Déploiement & Docker, Docker Compose & Environnements reproductibles \\
CI/CD & GitHub Actions, GHCR & Tests et publication des images \\
\bottomrule
\caption{Inventaire des technologies utilisées dans HARDTEC}
\end{longtable}

\chapter{Analyse fonctionnelle du système}
\section{Description du besoin métier}
Le besoin métier est double :
\begin{itemize}[leftmargin=1.3em]
    \item automatiser le tri et l’orientation des demandes de support ;
    \item augmenter la capacité d’analyse des cas historiques et améliorer les recommandations.
\end{itemize}
Une demande de support peut présenter une description textuelle libre. Cela rend le tri manuel complexe à cause de la variété des formulations, des synonymes, des erreurs de typographie et des niveaux d’urgence différents. Le système doit donc interpréter le texte de manière robuste.

\section{Flux métier actuel}
Le flux traditionnel suit cette logique :
\begin{enumerate}[leftmargin=1.8em]
    \item un utilisateur soumet un ticket ;
    \item l’opérateur lit la description ;
    \item le ticket est classé manuellement ;
    \item il est affecté à une équipe ;
    \item la solution est recherchée dans les cas précédents ;
    \item le traitement est suivi jusqu’à résolution.
\end{enumerate}
Le système proposé ajoute une couche d’assistance IA afin d’optimiser chaque étape.

\section{Flux métier cible}
Le flux cible devient :
\begin{enumerate}[leftmargin=1.8em]
    \item soumission du ticket ;
    \item classification automatique ;
    \item estimation de la priorité ;
    \item sélection de la file d’affectation ;
    \item recherche de cas historiques similaires ;
    \item calcul d’un risque d’incident ;
    \item génération d’une recommandation ;
    \item validation humaine avant décision finale.
\end{enumerate}

\chapter{Architecture générale de la plateforme}
\section{Vue d’ensemble}
L’architecture du projet est construite selon un modèle modulaire, avec une séparation claire entre :
\begin{itemize}[leftmargin=1.3em]
    \item couche interface utilisateur ;
    \item couche API métier ;
    \item couche intelligence artificielle ;
    \item couche données ;
    \item couche déploiement et supervision.
\end{itemize}

\begin{figure}[H]
\centering
\begin{tikzpicture}[node distance=1.0cm and 1.1cm]

\node[component] (ui) {Streamlit\\Dashboard};
\node[component,right=2.6cm of ui] (api) {FastAPI\\API};
\node[component,right=2.6cm of api] (agent) {Agent IA\\Orchestrateur};
\node[component,below=1.8cm of api] (ml) {ML Pipeline\\Type / Queue / Priority};
\node[component,below=1.8cm of agent] (rag) {RAG\\FAISS + embeddings};
\node[component,left=2.4cm of ml] (aiops) {AIOps\\Incident risk};
\node[data,below=2.0cm of ml] (sqlite) {SQLite\\Historique};
\node[data,below=2.0cm of rag] (mlflow) {MLflow\\Tracking};

\draw[arrow] (ui)--(api);
\draw[arrow] (api)--(agent);
\draw[arrow] (api)--(ml);
\draw[arrow] (agent)--(rag);
\draw[arrow] (agent)--(aiops);
\draw[arrow] (api)--(sqlite);
\draw[arrow] (ml)--(mlflow);
\end{tikzpicture}
\caption{Architecture logique de HARDTEC Intelligent Support \& AIOps}
\end{figure}

\section{Composants fonctionnels}
\subsection{Couche interface utilisateur}
La couche d’interface est assurée par Streamlit. Elle permet à l’utilisateur de saisir un ticket, de faire une prédiction et de consulter l’historique. Elle repose sur des requêtes HTTP vers l’API FastAPI.

\subsection{Couche API métier}
L’API FastAPI centralise les services de classification, d’historique, de statistiques et de recherche contextuelle. Elle expose des endpoints REST et est la principale interface de communication entre les composants applicatifs et la couche de décision.

\subsection{Couche intelligence artificielle}
Cette couche contient :
\begin{itemize}[leftmargin=1.3em]
    \item le moteur de prédiction de type de ticket ;
    \item le prédicteur de file d’affectation ;
    \item le prédicteur de priorité ;
    \item le moteur RAG basé sur FAISS ;
    \item le module AIOps d’évaluation du risque ;
    \item l’agent IA orchestrateur.
\end{itemize}

\subsection{Couche données}
Les données sont stockées localement dans SQLite pour la traçabilité des prédictions. Les fichiers de modèles et d’index vectoriels sont stockés dans le dossier \texttt{models}. MLflow conserve les métriques et les expérimentations associées au projet.

\chapter{Analyse technique détaillée du projet}
\section{Structure du dépôt}
L’application est structurée selon une organisation modulaire, notamment :
\begin{itemize}[leftmargin=1.3em]
    \item \texttt{src/api/} : API FastAPI ;
    \item \texttt{src/agent/} : service agent IA ;
    \item \texttt{src/pipeline/} : pipeline de prédiction ;
    \item \texttt{src/rag/} : moteur RAG et index vectoriel ;
    \item \texttt{src/aiops/} : détection d’anomalies et prédiction de risque ;
    \item \texttt{streamlit/} : interface utilisateur ;
    \item \texttt{data/} : données brutes, traitées et historiques ;
    \item \texttt{models/} : modèles entraînés ;
    \item \texttt{Docker/} : conteneurisation ;
    \item \texttt{.github/workflows/} : CI/CD ;
    \item \texttt{rapport\_pfa/} : version LaTeX du mémoire.
\end{itemize}

\section{API FastAPI}
Le fichier \texttt{src/api/main.py} constitue le cœur du backend. Il expose plusieurs endpoints essentiels :
\begin{itemize}[leftmargin=1.3em]
    \item \texttt{/} : information générale sur le service ;
    \item \texttt{/health} : vérification de l’état des composants ;
    \item \texttt{/predict} : prédiction d’un ticket ;
    \item \texttt{/history} : historique des prédictions ;
    \item \texttt{/stats} : statistiques agrégées ;
    \item \texttt{/rag/search} : recherche historique sémantique ;
    \item \texttt{/aiops/predict-risk} : prédiction de risque ;
    \item \texttt{/support/analyze} : analyse complète d’un ticket.
\end{itemize}

L’API est conçue pour être robuste :
\begin{itemize}[leftmargin=1.3em]
    \item validation avec Pydantic ;
    \item gestion explicite des exceptions ;
    \item stockage local SQLite pour les prédictions ;
    \item fallback sur composants optionnels ;
    \item retour structuré JSON pour chaque endpoint.
\end{itemize}

\begin{figure}[H]
\centering
\begin{tikzpicture}[node distance=1.1cm and 1.2cm]
\node[component] (client) {Client / UI};
\node[component,right=2.2cm of client] (api) {FastAPI};
\node[component,below=1.3cm of api] (ml) {Ticket ML};
\node[component,right=2.4cm of api] (db) {SQLite};
\node[component,below=1.3cm of db] (rag) {RAG};
\node[component,below=1.3cm of ml] (aiops) {AIOps};
\draw[arrow] (client)--(api);
\draw[arrow] (api)--(ml);
\draw[arrow] (api)--(db);
\draw[arrow] (api)--(rag);
\draw[arrow] (api)--(aiops);
\end{tikzpicture}
\caption{Interconnexion de la couche API}
\end{figure}

\section{Pipeline de prédiction}
Le moteur principal est situé dans \texttt{src/pipeline/predictor.py}. Il charge trois modèles :
\begin{itemize}[leftmargin=1.3em]
    \item modèle de type de ticket ;
    \item modèle de file d’affectation ;
    \item modèle de priorité.
\end{itemize}
La logique de prédiction suit l’ordre suivant :
\begin{enumerate}[leftmargin=1.8em]
    \item validation du texte fourni ;
    \item chargement des artefacts ;
    \item prédiction du type ;
    \item construction de métadonnées ;
    \item prédiction de la file ;
    \item prédiction de la priorité ;
    \item calcul de la confiance ;
    \item génération d’une recommandation.
\end{enumerate}

\section{Modèle de classification}
Le système repose sur des modèles supervisés issus de scikit-learn. Les tickets sont transformés en vecteurs TF-IDF puis utilisés pour entraîner les classifieurs. L’objectif est de capturer les représentations textuelles et de déterminer automatiquement la catégorie la mieux adaptée.

Le code présente une logique de pipeline où :
\begin{equation}
\text{Ticket} \rightarrow \text{TF-IDF} \rightarrow \text{Type} \rightarrow \text{Queue} \rightarrow \text{Priority}
\end{equation}

La sortie est structurée sous forme d’un dictionnaire contenant :
\begin{itemize}[leftmargin=1.3em]
    \item \texttt{ticket\_type} ;
    \item \texttt{queue} ;
    \item \texttt{priority} ;
    \item \texttt{confidence} ;
    \item \texttt{recommendation}.
\end{itemize}

\section{Moteur RAG}
Le composant RAG, dans \texttt{src/rag/rag\_engine.py}, est chargé de comparer la requête avec des tickets historiques. Il exploite :
\begin{itemize}[leftmargin=1.3em]
    \item un modèle d’embedding : \texttt{SentenceTransformer} ;
    \item un index vectoriel FAISS ;
    \item une base de documents historiques.
\end{itemize}

Le mécanisme fonctionne ainsi :
\begin{enumerate}[leftmargin=1.8em]
    \item encodage du ticket actuel ;
    \item recherche des meilleurs voisins dans l’index ;
    \item récupération des tickets textuellement proches ;
    \item extraction des métadonnées et solutions historiques ;
    \item renvoi des cas les plus utiles à l’agent.
\end{enumerate}

\begin{figure}[H]
\centering
\begin{tikzpicture}[node distance=1.1cm and 1.3cm]
\node[component] (query) {Ticket actuel};
\node[component,right=2.1cm of query] (embed) {Embedding};
\node[component,right=2.1cm of embed] (faiss) {Index FAISS};
\node[component,below=1.5cm of faiss] (docs) {Documents historiques};
\node[component,below=1.5cm of embed] (res) {Tickets similaires};
\draw[arrow] (query)--(embed);
\draw[arrow] (embed)--(faiss);
\draw[arrow] (faiss)--(docs);
\draw[arrow] (faiss)--(res);
\end{tikzpicture}
\caption{Architecture du moteur de recherche RAG}
\end{figure}

\section{AIOps et risque d’incident}
La logique AIOps complète le système en analysant les signaux de performance et d’incident. Le projet intègre des modules de détection d’anomalies et de risque. Cette couche permet d’anticiper les problèmes plutôt que de se limiter au traitement réactif. Les signaux peuvent être extraits de données temporelles, de métriques ou de corrélations entre événements.

\subsection{Pourquoi AIOps dans le projet ?}
L’objectif est de détecter les symptômes d’un incident avant qu’il ne devienne critique. Un ticket est alors enrichi par le niveau de risque du système. La solution permet de réduire les temps d’arrêt, d’améliorer la qualité de service et de supporter les équipes techniques dans les situations complexes.

\section{Stockage et traçabilité}
Le dossier \texttt{src/data/ticket\_store.py} met en place un stockage local SQLite, utilisé pour enregistrer les prédictions. L’intérêt est de conserver l’historique des tickets traités avec leurs prédictions, recommandations et timestamps. Cela permet de :
\begin{itemize}[leftmargin=1.3em]
    \item revoir les décisions précédentes ;
    \item mesurer la qualité des prédictions ;
    \item corriger les erreurs ;
    \item établir une base d’apprentissage et de validation.
\end{itemize}

\chapter{Agent IA intelligent}
\section{Rôle de l’agent}
Le service d’agent, dans \texttt{src/agent/main.py}, constitue une couche d’orchestration intelligente. Son rôle est de coordonner les différentes capacités du système sans les remplacer. Il complète l’architecture en agissant comme un assistant de décision pour le support.

L’agent reçoit un ticket, effectue la prédiction, récupère le contexte historique, construit un plan d’action et renvoie une réponse structurée. Cette réponse exploite, selon les cas :
\begin{itemize}[leftmargin=1.3em]
    \item les résultats du classifieur ;
    \item les tickets historiques similaires ;
    \item un résumé optionnel généré par un LLM externe ;
    \item des recommandations orientées opérationnelles.
\end{itemize}

\section{Architecture fonctionnelle de l’agent}
\begin{figure}[H]
\centering
\begin{tikzpicture}[node distance=1.15cm and 1.2cm]
\node[component] (user) {Ticket utilisateur};
\node[component,right=2.2cm of user] (agent) {Agent IA};
\node[component,right=2.2cm of agent] (ml) {Prédiction ML};
\node[component,below=1.5cm of ml] (rag) {RAG historique};
\node[component,below=1.5cm of agent] (plan) {Plan d’action};
\node[component,left=2.2cm of plan] (human) {Validation humaine};
\draw[arrow] (user)--(agent);
\draw[arrow] (agent)--(ml);
\draw[arrow] (agent)--(rag);
\draw[arrow] (agent)--(plan);
\draw[arrow] (human)--(agent);
\end{tikzpicture}
\caption{Flux de traitement de l’agent IA}
\end{figure}

\section{Plan d’action généré par l’agent}
Le plan d’action est construit sur la base de la priorité et de la file. Si le ticket est urgent, l’agent recommande une escalade immédiate. S’il est moins critique, il propose une prise en charge standard avec revue du contexte historique. Cette logique rend l’agent utile pour la priorisation opérationnelle et la meilleure allocation des ressources.

\section{Risque et supervision}
L’agent ne déclenche pas automatiquement des actions irréversibles. Il exige validation humaine. Cela correspond à une approche de sécurité et de responsabilité. En pratique, l’agent est un assistant de décision, pas une entité autonome de production.

\chapter{Interface utilisateur Streamlit}
\section{Objectif de l’interface}
L’interface \texttt{streamlit/app.py} est conçue pour la démonstration et la consultation. Elle permet de tester le système, d’analyser un ticket, d’afficher les résultats et d’explorer les performances. Elle sert à vulgariser l’usage de l’outil tout en restant compatible avec des besoins de support opérationnel.

\section{Fonctionnalités de l’interface}
\begin{itemize}[leftmargin=1.3em]
    \item saisie d’un ticket ;
    \item appel de l’API FastAPI ;
    \item affichage du type, de la file et de la priorité ;
    \item visualisation des statistiques et tendances ;
    \item affichage des tickets historiques similaires ;
    \item présentation du risque d’incident ;
    \item intégration d’un design simple et professionnel.
\end{itemize}

\section{Architecture front-end}
Le front-end interagit avec le backend par HTTP. La logique de communication s’effectue via la variable d’environnement \texttt{FASTAPI\_URL}. Cela permet de déployer l’interface et l’API indépendamment, tout en maintenant une communication stable.
\chapter{Base de données et stockage des prédictions}
\section{Captures d’écran de l’application}
Les captures suivantes documentent l’utilisation réelle de la plateforme. Le dépôt ne contient pas encore d’images exportées ; les emplacements ci-dessous sont prévus pour recevoir les screenshots pris pendant la démonstration. La macro \texttt{\\screenshot} affiche automatiquement l’image lorsqu’elle est déposée dans le dossier prévu.
\screenshot{streamlit-dashboard.png}{Vue générale du tableau de bord Streamlit avec saisie d’un ticket et affichage des résultats.}{Tableau de bord Streamlit}
\screenshot{prediction-result.png}{Résultat d’une prédiction : type, priorité, file, confiance et recommandation.}{Résultat de classification d’un ticket}
\screenshot{rag-results.png}{Liste des tickets historiques similaires retournés par le moteur RAG.}{Résultats de recherche sémantique}
\section{SQLite comme base locale}
Le projet utilise SQLite comme solution de stockage locale. C’est une approche légère, rapide et adaptée au contexte de démonstration et de validation. Elle permet de stocker les tickets prédits et les informations associées, sans dépendre d’un système de base de données centralisé.

\section{Structure de données}
Chaque prédiction est conservée avec :
\begin{itemize}[leftmargin=1.3em]
    \item texte du ticket ;
    \item type prédit ;
    \item file affectée ;
    \item priorité ;
    \item confiance ;
    \item recommandation ;
    \item source d’origine ;
    \item timestamp.
\end{itemize}

\section{Bénéfices}
Cette stratégie favorise :
\begin{itemize}[leftmargin=1.3em]
    \item le suivi d’historique ;
    \item la compréhension du comportement du système ;
    \item la génération de statistiques ;
    \item la préparation des données pour l’amélioration future.
\end{itemize}

\chapter{Suivi des expériences et reproductibilité}
\section{MLflow}
MLflow est utilisé pour le suivi des expériences et la gestion des artefacts. Il permet d’enregistrer les modèles, les métriques, les paramètres et les sorties d’entraînement. Cela garantit une traçabilité de la chaîne de développement et aide les équipes à comparer les performances entre versions.

\section{Intégration dans le projet}
Les modèles sont suivis avec leurs métriques telles que :
\begin{itemize}[leftmargin=1.3em]
    \item accuracy ;
    \item F1-score ;
    \item précision ;
    \item rappel ;
    \item paramètres d’entraînement.
\end{itemize}

\section{Impact pour le projet}
MLflow facilite :
\begin{itemize}[leftmargin=1.3em]
    \item le benchmarking des modèles ;
    \item le debugging des performances ;
    \item la comparaison des essais ;
    \item la reproductibilité des expériences ;
    \item la mise en production de versions validées.
\end{itemize}
\screenshot{mlflow-experiment.png}{Vue de l’expérience hardtec-ticketing, des runs et des métriques enregistrées.}{Suivi expérimental dans MLflow}

\chapter{Ingénierie des données et orchestration}
\section{Chaîne de données}
Le dépôt contient plusieurs niveaux de données : fichiers bruts, données nettoyées, tables analytiques, résultats de prédiction et sorties AIOps. Les fichiers CSV du dossier \texttt{data/processed/} représentent les résultats intermédiaires et finaux utilisés par les services et les analyses.

\begin{figure}[H]
\centering
\begin{tikzpicture}[node distance=1.0cm and 1.0cm]
\node[component] (raw) {Données brutes};
\node[component,right=1.8cm of raw] (dbt) {dbt\\Bronze / Silver / Gold};
\node[component,right=1.8cm of dbt] (ml) {Entraînement ML};
\node[component,below=1.5cm of ml] (aiops) {AIOps\\fenêtres et anomalies};
\node[component,below=1.5cm of dbt] (store) {CSV / SQLite};
\draw[arrow] (raw)--(dbt)--(ml);
\draw[arrow] (dbt)--(store);
\draw[arrow] (ml)--(aiops);
\end{tikzpicture}
\caption{Chaîne analytique et flux de données}
\end{figure}

\section{Modélisation dbt}
Le dossier \texttt{dbt/hardtec\_dbt/} contient des modèles SQL organisés en trois niveaux. Le niveau Bronze prépare les données sources, le niveau Silver applique les transformations et le niveau Gold produit des vues analytiques telles que \texttt{gold\_ticket\_analytics} et \texttt{gold\_ai\_predictions}. Le fichier \texttt{schema.yml} documente les modèles et constitue le point d’extension naturel pour ajouter des tests de qualité.

\section{Orchestration Airflow}
Le dossier \texttt{airflow/dags/} contient quatre DAGs spécialisés :
\begin{itemize}[leftmargin=1.3em]
    \item \texttt{hardtec\_ticket\_pipeline.py} pour la prédiction de tickets ;
    \item \texttt{hardtec\_rag\_pipeline.py} pour la construction et l’exploitation du contexte historique ;
    \item \texttt{hardtec\_aiops\_pipeline.py} pour les anomalies et les risques ;
    \item \texttt{hardtec\_complete\_pipeline.py} pour l’enchaînement complet.
\end{itemize}
Airflow permet de planifier les traitements, de visualiser leurs dépendances et de rejouer une étape en cas d’échec. Il constitue la couche d’orchestration batch du projet, complémentaire à FastAPI qui sert les prédictions en temps réel.

\chapter{AIOps, MicroSS et détection d’anomalies}
\section{Données MicroSS}
Le dossier \texttt{MicroSS/} contient des données de métriques sélectionnées et des résultats liés à l’analyse de services conteneurisés. Les métriques observées incluent notamment l’utilisation CPU et mémoire de services tels que Redis, webservice, logservice et dbservice. Elles alimentent l’analyse temporelle et la construction de signaux de risque.

\section{Détection temporelle}
Le dossier \texttt{metric\_detection/} regroupe des séries de référence correspondant à différents comportements : données périodiques, linéaires, en escalier, partiellement stationnaires et présentant un faible rapport signal/bruit. Ces données servent à comparer ou à valider les méthodes de détection d’anomalies dans plusieurs régimes temporels.

\section{Résultats AIOps}
Les scripts du dossier \texttt{src/aiops/} construisent des fenêtres temporelles, détectent les anomalies, agrègent les signaux et produisent des prédictions de risque. Les sorties sont enregistrées dans \texttt{data/processed/}, notamment sous forme d’incidents, anomalies, fenêtres et scores de risque. Cette chaîne permet de relier les observations techniques à la décision de support.

\begin{figure}[H]
\centering
\begin{tikzpicture}[node distance=1.0cm and 1.0cm]
\node[component] (metrics) {Métriques MicroSS};
\node[component,right=1.8cm of metrics] (windows) {Fenêtres temporelles};
\node[component,right=1.8cm of windows] (anom) {Anomalies};
\node[component,below=1.5cm of anom] (risk) {Risque d’incident};
\node[component,below=1.5cm of windows] (inc) {Incidents corrélés};
\draw[arrow] (metrics)--(windows)--(anom)--(risk);
\draw[arrow] (windows)--(inc)--(risk);
\end{tikzpicture}
\caption{Pipeline AIOps de la métrique au risque d’incident}
\end{figure}

\chapter{Scripts, logs et artefacts}
\section{Scripts d’analyse}
Les fichiers Python situés à la racine et dans \texttt{scripts/} servent à inspecter les données, analyser les blocs actifs, calculer des statistiques et préparer les sorties analytiques. Ils constituent des outils d’exploration et de contrôle, distincts des services applicatifs.

\section{Logs et résultats}
Les dossiers \texttt{logs/}, \texttt{log/}, \texttt{reports/} et \texttt{aiops\_analysis/} regroupent les traces, résultats intermédiaires et synthèses d’analyse. Leur présence complète la chaîne de traçabilité : données d’entrée, transformation, modèle, prédiction, contrôle et restitution.

\section{Limites de l’inventaire}
Les fichiers générés, les caches, les environnements virtuels et les volumes techniques ne sont pas décrits comme des composants fonctionnels. Ils sont nécessaires à l’exécution locale, mais ne constituent pas des briques de l’architecture métier.

\chapter{Conteneurisation et déploiement Docker}
\section{Objectif}
Le projet a été pensé pour être exécutable et déployable dans un environnement standardisé. La conteneurisation permet de gérer le projet avec moins d’écarts entre les environnements de développement, de test et de production.

\section{Services conteneurisés}
Le fichier \texttt{Docker/docker-compose.yml} définit trois services principaux :
\begin{itemize}[leftmargin=1.3em]
    \item API FastAPI sur le port 8000 ;
    \item agent IA sur le port 8001 ;
    \item MLflow sur le port 5000.
\end{itemize}

\begin{figure}[H]
\centering
\begin{tikzpicture}[node distance=1.0cm and 1.2cm]
\node[component] (api) {API\8000};
\node[component,right=2.5cm of api] (agent) {Agent\8001};
\node[component,right=2.5cm of agent] (mlflow) {MLflow\5000};
\node[data,below=1.7cm of agent] (models) {Modèles};
\node[data,below=1.7cm of mlflow] (vol) {Volume MLflow};
\draw[arrow] (api)--(agent);
\draw[arrow] (agent)--(mlflow);
\draw[arrow] (models)--(api);
\draw[arrow] (models)--(agent);
\draw[arrow] (vol)--(mlflow);
\end{tikzpicture}
\caption{Déploiement Docker de la plateforme}
\end{figure}

\section{Justification du choix de l’architecture}
Le système ne nécessite pas MinIO dans sa première version. Dans cette architecture, les données structurées et historiques restent hébergées dans le dépôt du projet, la base SQLite et les artefacts MLflow. Le principe est de garder une solution légère, simple et facilement testable. Cela correspond bien au contexte de recherche et de démonstration en stage.

\section{Commande de lancement}
\begin{lstlisting}[language=bash]
docker compose -f Docker/docker-compose.yml up --build
\end{lstlisting}
Ce lancement crée les services de l’API, de l’agent et de l’observabilité MLflow.

\chapter{CI/CD et automatisation}
\section{Objectifs de l’automatisation}
Le workflow GitHub Actions est défini dans \texttt{.github/workflows/ci-cd.yml}. Il a pour but de vérifier automatiquement la qualité du code, la conformité de l’architecture et la capacité de déploiement du projet.

\section{Étapes de validation}
Le pipeline contient les étapes suivantes :
\begin{enumerate}[leftmargin=1.8em]
    \item téléchargement du code ;
    \item installation des dépendances ;
    \item compilation Python ;
    \item exécution des tests ;
    \item validation de la configuration Docker Compose ;
    \item construction et publication de l’image sur GHCR.
\end{enumerate}

\begin{figure}[H]
\centering
\begin{tikzpicture}[node distance=1.0cm and 1.1cm]
\node[component] (push) {Push ou PR};
\node[component,right=2.5cm of push] (tests) {Tests\+ compile};
\node[component,right=2.5cm of tests] (docker) {Compose\validation};
\node[component,right=2.5cm of docker] (build) {Build image};
\node[component,below=1.7cm of build] (ghcr) {GitHub Container\Registry};
\draw[arrow] (push)--(tests)--(docker)--(build);
\draw[arrow] (build)--(ghcr);
\end{tikzpicture}
\caption{Chaîne CI/CD mise en place pour le projet}
\end{figure}

\section{Bénéfices}
Cette automatisation réduit les risques de régression, permet un contrôle qualité systématique et prépare le projet à une mise en production plus maîtrisée.

\chapter{Tests et validation du système}
\section{Types de validation}
Le système a été validé à plusieurs niveaux :
\begin{itemize}[leftmargin=1.3em]
    \item validation syntaxique ;
    \item validation des endpoints API ;
    \item validation de la configuration Docker ;
    \item validation du moteur de prédiction ;
    \item validation du service agent.
\end{itemize}

\section{Exemples de tests}
Les tests du dépôt couvrent :
\begin{itemize}[leftmargin=1.3em]
    \item stockage local ;
    \item prédicteur de tickets ;
    \item configuration ;
    \item endpoints API ;
    \item disponibilité du service agent.
\end{itemize}

\section{Interprétation des résultats}
La validation montre que le système est capable de démarrer et de produire des prédictions cohérentes sur des tickets de support IT. Les traces MLflow permettent de vérifier la présence des expériences, des métriques et des exécutions de modèles. La présence d’un historique de prédictions renforce la capacité de monitoring et d’amélioration continue.

\chapter{Analyse critique et limites}
\section{Points forts}
\begin{itemize}[leftmargin=1.3em]
    \item architecture modulaire et lisible ;
    \item intégration de plusieurs briques IA ;
    \item API robuste et testable ;
    \item présence d’un moteur d’assistance IA ;
    \item traçabilité via MLflow ;
    \item déploiement reproductible avec Docker ;
    \item automatisation via GitHub Actions.
\end{itemize}

\section{Limites actuelles}
\begin{itemize}[leftmargin=1.3em]
    \item dépendance aux modèles et aux artefacts disponibles localement ;
    \item besoin de données historiques de meilleure qualité pour enrichir le RAG ;
    \item nécessité de sécuriser les variables d’environnement ;
    \item besoin d’une authentification plus stricte pour l’API en production ;
    \item nécessité d’un monitoring avancé (logs, métriques, alertes).
\end{itemize}

\section{Perspectives d’évolution}
Le projet peut évoluer vers :
\begin{itemize}[leftmargin=1.3em]
    \item l’intégration d’authentification OAuth2 ou JWT ;
    \item un panel de monitoring avec Prometheus et Grafana ;
    \item un réseau plus complet de données structurant les incidents ;
    \item un mécanisme de revue par l’équipe support ;
    \item une version de production orientée Kubernetes ou Container Apps ;
    \item une meilleure gouvernance des modèles et des validations.
\end{itemize}

\chapter{Conclusion générale}
Le projet HARDTEC Intelligent Support \& AIOps se présente comme une solution complète d’assistance au support IT basée sur l’intelligence artificielle. Il couvre plusieurs dimensions : classification des demandes, prédiction de leur priorité, routage vers la bonne équipe, recherche de solutions historiques, analyse de métriques, détection d’anomalies, prédiction du risque et assistance intelligente. La plateforme est organisée selon une architecture modulaire et exploitable, avec des points forts clairs : simplicité d’intégration, traçabilité, reproductibilité et évolutivité.

Le travail réalisé lors de ce stage montre qu’il est possible de transformer un processus de support manuel en un environnement assisté par des technologies IA. L’architecture mise en place constitue une base solide pour des améliorations futures dans le domaine de l’AIOps, de la gestion intelligente des incidents et de la décision assistée par intelligence artificielle.

\chapter*{Conclusion personnelle}
\addcontentsline{toc}{chapter}{Conclusion personnelle}
Ce stage m’a permis de comprendre le rôle stratégique de l’intelligence artificielle dans les systèmes de support informatique. J’ai pu développer une solution de bout en bout, du modèle de classification jusqu’à l’interface utilisateur et au déploiement conteneurisé. Ce projet a renforcé mes compétences en développement Python, architecture logicielle, modélisation ML, intégration d’API et gestion de projet. Il m’a aussi permis de mieux comprendre les enjeux industriels et académiques liés à l’IA appliquée aux services digitaux.

\chapter*{Bibliographie}
\addcontentsline{toc}{chapter}{Bibliographie}
\begin{itemize}[leftmargin=1.3em]
    \item FastAPI Documentation, \\https://fastapi.tiangolo.com/
    \item Streamlit Documentation, \\https://streamlit.io/
    \item MLflow Documentation, \\https://mlflow.org/
    \item FAISS Documentation, \\https://github.com/facebookresearch/faiss
    \item SentenceTransformers Documentation, \\https://www.sbert.net/
    \item scikit-learn Documentation, \\https://scikit-learn.org/
    \item Docker Documentation, \\https://docs.docker.com/
    \item GitHub Actions Documentation, \\https://docs.github.com/actions
    \item HARDTEC MAROC, site officiel, \\https://hardtec.net
\end{itemize}

\chapter*{Annexes}
\addcontentsline{toc}{chapter}{Annexes}
\section*{Annexe A : Exemple de ticket traité par le système}
\begin{lstlisting}[language=text]
Ticket: Le serveur applicatif est lent et plusieurs utilisateurs
signalent des erreurs de connexion. La session est interrompue
au bout de quelques minutes. \n
Priorite predite: HIGH
File predite: Infrastructure
Type predite: Incident
\end{lstlisting}

\section*{Annexe B : Exemple de réponse de l’API}
\begin{lstlisting}[language=json]
{
  "ticket_text": "Le VPN est indisponible pour plusieurs utilisateurs",
  "type": "Incident",
  "priority": "HIGH",
  "queue": "Infrastructure",
  "confidence": 0.93,
  "recommendation": "Assign the ticket to the Infrastructure team with HIGH priority and monitor the incident through its resolution lifecycle."
}
\end{lstlisting}

\section*{Annexe C : Exemple de requête Agent IA}
\begin{lstlisting}[language=python]
POST http://localhost:8001/agent/analyze
{
  "ticket_text": "Le VPN est indisponible pour plusieurs utilisateurs",
  "top_k": 3,
  "allow_external_llm": false
}
\end{lstlisting}

\section*{Annexe D : Commandes de démonstration}
\begin{lstlisting}[language=bash]
# Démarrer la plateforme
cd hardtec-intelligent-ticketing
docker compose -f Docker/docker-compose.yml up --build

# Vérification API
curl http://localhost:8000/health

# Vérification agent
curl http://localhost:8001/health

# Vérification MLflow
http://localhost:5000
\end{lstlisting}

\end{document}
